\section{Workload Model}

In the present section, the workload model of Airbnb will be described, giving the different parameters with 
its corresponding value. Each value is justified with its proper reference, the data that wasn't publicly available
was estimated taking into account other parameters and data, its justification will be found under the corresponding
table.

% -----TRAFFIC AND LOAD-----
\subsection{Traffic and load}

The following table shows the value of the traffic that Airbnb experiments.

\begin{longtable}{|p{0.27\textwidth}|p{0.55\textwidth}|}
    \hline
    \textbf{Parameter} & 
    \textbf{Value} \\ 
    \hline
    \endfirsthead%

    \hline
    \textbf{Parameter} & 
    \textbf{Value} \\ 
    \hline
    \endhead%

    Average requests/hour &
    $\approx 934{,}000$ page requests/hour \cite{semrush_airbnb_2026}\\
    \hline

    Peak requests/hour &
    $\approx 9.3$ million page requests/hour \cite{semrush_airbnb_2026, airbnb_engineering_traffic_spikes} \\
    \hline

    Peak duration &
    $\approx 5$ minutes \cite{airbnb_engineering_traffic_spikes} \\
    \hline

    Burst rate &
    $10\times$ normal; $20\times$ stress case \cite{airbnb_engineering_traffic_spikes} \\
    \hline

    Seasonal traffic multiplier &
    $\approx 1.1\times$ \cite{airbnb_10k_2025} \\
    \hline

\end{longtable}

The value of average requests/hour was estimated using Semrush, which reports 93.03M monthly visits and 7.23 pages/visit. 
$93.03M \times 7.23 / 720 \approx 934{,}000$ page requests/hour. This represents observable web page 
traffic, not Airbnb's complete internal service-to-service traffic \cite{semrush_airbnb_2026}.

The peak requests/hour were estimated. Taking into account that the average web workload above is multiplied by a 
conservative $10\times$ burst factor. Airbnb engineering has documented traffic spikes of approximately 
$10$--$20\times$ during severe events. $934{,}000 \times 10 \approx 9.34$ million requests/hour 
\cite{semrush_airbnb_2026,airbnb_engineering_traffic_spikes}.

The burst rate was derived from Airbnb engineering, which has documented QPS increases in the $10$--$20\times$ range. 
$10\times$ is used as the normal capacity-planning burst and $20\times$ as an extreme stress scenario.
Source: Airbnb Engineering \cite{airbnb_engineering_traffic_spikes}.

The seasonal traffic multiplier was estimated. Airbnb states that its business is seasonal and that Nights and 
Seats Booked are generally highest in Q1--Q3, with Q3 being the peak travel season in North America and EMEA. 
A $1.1\times$ multiplier provides a conservative approximation of normal seasonal variation without treating the 
annual peak as an outage-level burst \cite{airbnb_10k_2025}.

% -----USER ACTIVITY-----
\subsection{User Activity}

The following table shows the value of the user activity in the application.

\begin{longtable}{|p{0.27\textwidth}|p{0.55\textwidth}|}
    \hline
    \textbf{Parameter} & 
    \textbf{Value} \\
    \hline
    \endfirsthead

    \hline
    \textbf{Parameter} & 
    \textbf{Value} \\
    \hline
    \endhead

    Average session duration & 
    10 min 42 s \cite{semrush_airbnb_2026} \\ 
    \hline 
    
    Requests per session & 
    7.23 page requests/session \cite{semrush_airbnb_2026} \\ 
    \hline 
    
    Concurrent users (average) & 
    $\approx 23{,}000$ web sessions \cite{semrush_airbnb_2026} \\ 
    \hline 
    
    Concurrent users (peak) & 
    $\approx 60{,}000$ web sessions \cite{semrush_airbnb_2026} \\ 
    \hline 
    
\end{longtable}

The average session duration was recorded from a third-party measurement. Semrush reports an average visit 
duration of 10:42 for airbnb.com in August 2026 \cite{semrush_airbnb_2026}. 

The amount of requests per session was used from a third-party measurement. Semrush reports 7.23 pages per visit. 
This is interpreted as user-facing page requests rather than internal microservice requests \cite{semrush_airbnb_2026}.

The concurrent users were estimated, using approximately 93.03M monthly visits and an average session duration of 
10.7 minutes: \[ \frac{93.03M}{30.4\ days}\times\frac{10.7\ min}{1440\ min/day} \approx 22{,}800 \] 
This estimates simultaneously active web sessions and excludes users interacting exclusively through native applications.
Source: Semrush traffic and session-duration data \cite{semrush_airbnb_2026}.

The peak of the concurrent users was estimated. A peak concurrency factor of approximately $2.5\times$ the average 
is used: $22{,}800\times2.5\approx57{,}000$, rounded to approximately 60,000. The factor is intentionally lower than 
the $10\times$ request burst because more requests per user can occur without a proportional increase in active users
\cite{semrush_airbnb_2026}.


% -----Read/Write ratio----- 
\subsection{Read and Write Ratio}

The following table shows the value of the user read and write ratio of the application.

\begin{longtable}{|p{0.27\textwidth}|p{0.55\textwidth}|}
    \hline
    \textbf{Parameter} & 
    \textbf{Value} \\
    \hline
    \endfirsthead

    \hline
    \textbf{Parameter} & 
    \textbf{Value} \\
    \hline
    \endhead

    Read/write ratio & 
    $\approx 95:5$ \cite{airbnb_engineering_mussel} \\ 
    \hline 
    
    Request distribution by endpoint & 
    Search/browse 60\%; listing/details 20\%; 
    account/auth 8\%; booking/payment 5\%; 
    messaging 4\%; other 3\%~\cite{educative_airbnb_system_design} \\ 
    \hline 
    
    Cache hit ratio & 
    $\approx 85\%$ of cacheable reads~\cite{educative_airbnb_system_design} \\ 
    \hline

    
\end{longtable}

The read and write ratio was estimated from Airbnb internal infrastructure. Airbnb's Mussel key-value store reports 
average read QPS above 800k and write QPS above 35k, corresponding to approximately 96:4. A rounded 95:5 ratio 
is used for the general workload model because search and discovery are substantially more read-intensive than 
booking and data modification \cite{airbnb_engineering_mussel}. 

The request distribution was estimated. The distribution emphasizes search because Airbnb system-design and engineering
material identify search/discovery as the dominant read-heavy workload, while booking/payment represents a smaller 
but consistency-critical workload. Percentages are architectural assumptions for capacity modeling 
\cite{educative_airbnb_system_design}.

The cache hit ratio was estimated. Airbnb publicly documents extensive caching and derived-data infrastructure, 
but does not publish a single platform-wide cache hit ratio. An 85\% hit ratio is a reasonable architectural assumption 
for frequently accessed listing/search data and should be treated as a tunable model parameter 
\cite{educative_airbnb_system_design}.


% -----PERFORMANCE AND RELIABILITY----- 
\subsection{Performance and reliability}

The following table shows the value of the performance of the application.

\begin{longtable}{|p{0.27\textwidth}|p{0.55\textwidth}|}
    \hline
    \textbf{Parameter} & 
    \textbf{Value} \\
    \hline
    \endfirsthead

    \hline
    \textbf{Parameter} & 
    \textbf{Value} \\
    \hline
    \endhead

    Availability & 
    $\geq 99.9\%$ \cite{airbnb_engineering_mussel} \\ 
    \hline 
    
    Response time (p95) & 
    $\approx 1.5$ s user-facing target \cite{airbnb_engineering_mussel,educative_airbnb_system_design} \\ 
    \hline 
    
\end{longtable}

The availability was used from a report for a major Airbnb infrastructure component. Airbnb's Mussel service 
reports greater than 99.9\% availability. This value is used as a conservative platform workload target, 
but cannot be described as an officially published end-to-end Airbnb SLA \cite{airbnb_engineering_mussel}.

The response time, within a 95 percentile was estimated. Even though Airbnb does not publish a global end-to-end 
p95. Mussel reports average p95 read latency below 8 ms, demonstrating backend-storage performance, while an end-user 
request traverses multiple services, networks, caches and rendering layers. Therefore, approximately 1.5 s is used 
as a user-facing workload target rather than incorrectly treating 8 ms as total request latency
\cite{airbnb_engineering_mussel,educative_airbnb_system_design}.

% -----DATA----- 
\subsection{Data}

The following table shows the value of the data growth in the application.

\begin{longtable}{|p{0.27\textwidth}|p{0.55\textwidth}|}
    \hline
    \textbf{Parameter} & 
    \textbf{Value} \\
    \hline
    \endfirsthead

    \hline
    \textbf{Parameter} & 
    \textbf{Value} \\
    \hline
    \endhead
    
    Data growth & 
    $\approx 8$--$10\%$/year \cite{airbnb_10k_2025} \\ 
    \hline 
    
\end{longtable}

The data growth was estimated from reported business growth. Airbnb reported 533M Nights and Seats Booked in 
2025 versus 492M in 2024, an 8\% increase. A baseline of 8--10\% annual growth is therefore used for workload/data 
planning, recognizing that storage growth will not necessarily equal booking growth \cite{airbnb_10k_2025}.


% -----MEDIA----- 
\subsection{Media}

The following table shows the value of the information related to the media files in the application.

\begin{longtable}{|p{0.27\textwidth}|p{0.55\textwidth}|}
    \hline
    \textbf{Parameter} & 
    \textbf{Value} \\
    \hline
    \endfirsthead

    \hline
    \textbf{Parameter} & 
    \textbf{Value} \\
    \hline
    \endhead
    
    Maximum payload size & 
    10 MB/photo; 100 MB/video \cite{airbnb_resource_center_photos, airbnb_help_video_verification}\\ 
    \hline 
    
    File upload frequency & 
    $\approx 9{,}000$ photos/hour \cite{searchlogistics_airbnb_statistics, airbnb_resource_center_photos}\\ 
    \hline 
    
\end{longtable}

The file upload frequency was estimated. Considering Airbnb has more than 9M active listings according to current 
third-party industry data. Assuming roughly one photo update per listing per year and several photos per update produces 
an order of magnitude of several thousand photo uploads/hour; 9,000/hour is used as a conservative media-workload estimate
\cite{searchlogistics_airbnb_statistics, airbnb_resource_center_photos}.

% -----SEARCH----- 
\subsection{Search}

The following table shows the value of the searches in the application, and their complexity.

\begin{longtable}{|p{0.27\textwidth}|p{0.55\textwidth}|}
    \hline
    \textbf{Parameter} & 
    \textbf{Value} \\
    \hline
    \endfirsthead

    \hline
    \textbf{Parameter} & 
    \textbf{Value} \\
    \hline
    \endhead
    
    Search complexity & 
    High; geo-search + dates + availability + price + filters + ranking/personalization
    \cite{educative_airbnb_system_design} \\ 
    \hline
    
\end{longtable}

The search complexity was oriented as qualitative analysus, due to the many procedures it requires. Airbnb search must 
combine geographic location, dates, availability, pricing, amenities and ranking/personalization. This is therefore 
a multi-dimensional read-heavy workload requiring indexes, caching and precomputed/derived data 
\cite{educative_airbnb_system_design}. 


% -----GEOLOCATION----- 
\subsection{Geography}

The following table shows the value of the geography scope in the application.

\begin{longtable}{|p{0.27\textwidth}|p{0.55\textwidth}|}
    \hline
    \textbf{Parameter} & 
    \textbf{Value} \\
    \hline
    \endfirsthead

    \hline
    \textbf{Parameter} & 
    \textbf{Value} \\
    \hline
    \endhead
    
    Geographic scope & 
    More than 220 countries and regions \cite{airbnb_10k_2025} \\ 
    \hline 
    
    Traffic by region & 
    North America 30\%; EMEA 40\%; Latin America 17\%; Asia Pacific 13\% \cite{airbnb_10k_2025} \\ 
    \hline
    
\end{longtable}

The traffic by region was used from a reported business-workload proxy. Airbnb reports 2025 Nights and Seats 
Booked by host-listing region: North America 158M (30\%), EMEA 215M (40\%), Latin America 90M (17\%), and 
Asia Pacific 70M (13\%). These figures describe booking activity rather than HTTP traffic, so they should be identified
as a traffic-distribution proxy \cite{airbnb_10k_2025}.